\reading{CR-21}{The Evolution of Stress Testing Counterparty Exposures}{DEFER}{3}

\begin{crkeybox}[Learning objectives for this reading]
\begin{enumerate}[label=\textbf{\alph*.},leftmargin=2em]
\item Differentiate among current exposure, peak exposure, expected exposure, and expected positive exposure.
\item Explain the treatment of counterparty credit risk (CCR) both as a credit risk and as a market risk and describe its implications for trading activities and risk management for a financial institution.
\item Describe a stress test that can be performed on a loan portfolio, and on a derivative portfolio.
\item Differentiate between stressed expected loss and stress loss of a credit portfolio, and calculate the stress loss on a loan portfolio and the stress loss on a derivative portfolio.
\item Describe a stress test that can be performed on CVA.
\item Calculate the stressed CVA and the stress loss on CVA.
\item Calculate the DVA and explain how stressing DVA enters into aggregating stress tests of CCR.
\item Describe the common pitfalls in stress testing CCR.
\end{enumerate}
\end{crkeybox}

\begin{crtrapbox}[Single-layer reading]
CR-21 exists only in the Crash layer.
\end{crtrapbox}

% =====================================================================
\lo{CR-21 a}{Differentiate among current exposure, peak exposure, expected exposure, and expected positive exposure.}

\begin{center}
\footnotesize
\begin{tabular}{L{3.4cm} L{11.1cm}}
\toprule
\thead{Measure} & \thead{Definition} \\
\midrule
\term{Current exposure} & Maximum potential loss upon \emph{immediate} default:
  zero, or the market value of the transaction \\
\term{Peak exposure} & Potential exposure at a \emph{high percentile}, for
  example 95\%, over a future timeframe \\
\term{Expected exposure} & \emph{Average} potential exposure over a future
  timeframe \\
\term{Expected positive exposure (EPE)} & Weighted average of expected exposures,
  often used for \emph{capital requirements} \\
\bottomrule
\end{tabular}
\end{center}
\figcap{Four measures ordered by what they average over. Current exposure is a
point in time with no averaging at all; peak is a percentile of the future
distribution; expected exposure is its mean at a date; EPE averages that mean
across dates. The fuller treatment of these metrics is at \textbf{CR-19 a}.}

% =====================================================================
\lo{CR-21 b}{Explain the treatment of counterparty credit risk (CCR) both as a credit risk and as a market risk and describe its implications for trading activities and risk management for a financial institution.}

Credit valuation adjustment represents the \term{market value} of counterparty
credit risk. A financial institution could view CCR as either credit risk or
market risk, although it should consider both.

\begin{center}
\footnotesize
\begin{tabular}{L{3.0cm} L{11.5cm}}
\toprule
\thead{Treated as} & \thead{Implication} \\
\midrule
\term{Credit risk} & Exposes the institution to changes in CVA. CVA should
  therefore be included in valuing a derivatives portfolio; otherwise the
  portfolio could experience large changes in market value \\
\term{Market risk} & Allows the institution to \emph{hedge} market risk losses,
  but leaves it exposed to declines in counterparty creditworthiness and to
  default \\
\term{Both} & Prudent, but complex and difficult to interpret \\
\bottomrule
\end{tabular}
\end{center}
\figcap{Each single-lens treatment leaves the other exposure uncovered, and the
two-lens treatment buys coverage at the price of interpretability. There is no
clean answer, which is itself the examinable point.}

% =====================================================================
\lo{CR-21 c}{Describe a stress test that can be performed on a loan portfolio, and on a derivative portfolio.}

\begin{enumerate}[label=\alph*)]
\item Financial institutions conduct stress tests on current exposures by
adjusting counterparty portfolios under scenarios like equity crashes, to
identify those with the highest risk.
\item These results, summarized in detailed tables, inform management about which
counterparties are most vulnerable to specific market shocks.
\item However, these stress tests often struggle with \term{result aggregation}
and fail to account for counterparty \term{credit quality} and \term{wrong-way
risk}, making it challenging to gauge the true risk from these tests alone.
\end{enumerate}

% =====================================================================
\lo{CR-21 d}{Differentiate between stressed expected loss and stress loss of a credit portfolio, and calculate the stress loss on a loan portfolio and the stress loss on a derivative portfolio.}

\begin{crdefbox}[The distinction the objective turns on]
\term{Stressed expected loss} $(\mathrm{EL}_S)$ is the expected loss recomputed
under the stressed inputs. It is a \emph{level}.

\smallskip
\term{Stress loss} is the \emph{difference}, $\mathrm{EL}_S - \mathrm{EL}$: how
much worse the stressed world is than the current one. It is a \emph{change}.
\end{crdefbox}

\subsection*{Stress test on a loan portfolio}

A stress test could take exposure at default and loss given default as
\emph{deterministic} and focus on stresses where the probability of default is
subject to a stress. The probability of default is then taken to be a function of
other variables, for example the exchange rate or the unemployment rate.

\begin{crfmlbox}[Stressed expected loss and stress loss, loan portfolio]
\[ \mathrm{EL}_S \;=\; \sum_{i=1}^{N} \mathrm{PD}_i^{S}\times \mathrm{EAD}_i
   \times \mathrm{LGD}_i \]
\[ \text{stress loss} \;=\; \mathrm{EL}_S - \mathrm{EL} \]
\vk{The superscript $S$ marks the stressed quantity. Only PD carries the
superscript here, because EAD and LGD were assumed deterministic.}
\end{crfmlbox}

\subsection*{Stress test on a derivative portfolio}

Here exposure at default is \emph{stochastic} and depends on the levels of market
variables, so the expected loss is built from EPE scaled by an \term{alpha
factor}. The institution can stress the PD as in the loan case, and can
additionally stress the EPE.

\begin{crfmlbox}[Expected loss and stressed expected loss, derivative portfolio]
\[ \mathrm{EL} \;=\; \sum_{i=1}^{N} \mathrm{PD}_i \times
   \bigl(\mathrm{EPE}_i \times \alpha\bigr)\times \mathrm{LGD}_i \]
\[ \mathrm{EL}_S \;=\; \sum_{i=1}^{N} \mathrm{PD}_i^{S} \times
   \bigl(\mathrm{EPE}_i^{S} \times \alpha\bigr)\times \mathrm{LGD}_i \]
\vk{$\alpha$ is the multiplier that converts effective EPE into an exposure at
default. Note that \emph{two} quantities carry the stress superscript in the
derivative case, PD and EPE, against only one in the loan case. The stress loss is
then computed the same way, $\mathrm{EL}_S - \mathrm{EL}$.}
\end{crfmlbox}

\begin{crtrapbox}[Which inputs are stressed is the discriminating fact]
Loan portfolio: stress PD only; EAD and LGD are deterministic.

\smallskip
Derivative portfolio: stress PD \emph{and} EPE, because exposure at default is
itself stochastic and depends on market variables.
\end{crtrapbox}

% =====================================================================
\lo{CR-21 e}{Describe a stress test that can be performed on CVA.}

\begin{enumerate}
\item First, use a common simplified formula for CVA to a counterparty that omits
wrong-way risk.
\item Second, aggregate across $N$ counterparties.
\item Third, calculate stressed CVA by applying an \term{instantaneous shock} to
some of the market variables. The stress could affect \emph{EE} and \emph{PD}.
\item Fourth, compute the stress loss on CVA.
\end{enumerate}

% =====================================================================
\lo{CR-21 f}{Calculate the stressed CVA and the stress loss on CVA.}

\begin{crfmlbox}[Stressed CVA aggregated across counterparties]
\[ \mathrm{CVA}_S \;=\; -\sum_{n=1}^{N} \mathrm{LGD}^{*}_{n}\times
   \sum_{j=1}^{T} \mathrm{EPE}^{S}_{n}(t_j)\times
   \mathrm{PD}^{S}_{n}(t_{j-1},\,t_j) \]
\[ \text{stress loss on CVA} \;=\; \mathrm{CVA}_S - \mathrm{CVA} \]
\vk{$\mathrm{LGD}^{*}_{n}$ is the loss given default for counterparty $n$. The
inner sum runs over time buckets $j$ and the outer sum over counterparties $n$.
Both EPE and PD carry the stress superscript, matching the derivative-portfolio
case at objective d, and the leading minus sign is the same convention as the
unstressed CVA at \textbf{CR-20 c}.}
\end{crfmlbox}

% =====================================================================
\lo{CR-21 g}{Calculate the DVA and explain how stressing DVA enters into aggregating stress tests of CCR.}

The DVA calculation itself is set out at \textbf{CR-20 e and g}, where
$\mathrm{DVA} = -\mathrm{LGD}_I \sum \mathrm{ENE}(t_i)\times
\mathrm{PD}_I(t_{i-1},t_i)$ and comes out positive because ENE is negative.

\begin{crkeybox}[How stressed DVA enters the aggregation]
Stress loss on DVA can be calculated in a similar way to CVA losses: compute a
stressed BCVA and subtract the current BCVA.

\smallskip
\[ \text{stress loss on BCVA} \;=\; \mathrm{BCVA}_S - \mathrm{BCVA} \]

Because $\mathrm{BCVA} = \mathrm{CVA} + \mathrm{DVA}$ and the two terms carry
opposite signs, a stress that widens \emph{both} the counterparty's and the
institution's own credit spreads produces offsetting effects. This is why DVA
cannot simply be ignored when aggregating CCR stress results: dropping it
overstates the stress loss.
\end{crkeybox}

% =====================================================================
\lo{CR-21 h}{Describe the common pitfalls in stress testing CCR.}

\begin{enumerate}
\item CCR is often \term{not integrated} with stress tests for loan portfolios or
trading positions.
\item Institutions commonly use \term{current exposure} instead of expected
exposure or expected positive exposure, which can result in significant errors.
\item Relying on current exposure is especially problematic for
\term{at-the-money} derivative positions.
\item The use of \term{delta sensitivities} in exposure calculations poses
challenges due to their nonlinearity; linearizing these sensitivities can lead to
considerable inaccuracies.
\end{enumerate}

\begin{crtrapbox}[Why at-the-money is the worst case for current exposure]
An at-the-money position has a current exposure close to zero and a potential
future exposure that is not. Using current exposure therefore reports the
position as harmless precisely where the future uncertainty is greatest, which is
pitfall 2 and pitfall 3 acting together.
\end{crtrapbox}
